\naslov{Simetrični problem lastnih vrednosti}

Matrika $A = A^T$ ima same realne lastne vrednosti, torej jih lahko uredimo kot
$\lambda_n \le \cdots \le \lambda_1$.
Lastne vektorje lahko izberemo tako, da so ortonormirani, $A x_i = \lambda_i
x_i$, $x_i^T x_j = \delta_{ij}$.

\begin{lema}
  Za Rayleighov kvocient velja $\lambda_n \le \rho(x, A) \le \lambda_1$.
\end{lema}

\begin{proof}
  Zapišemo $x$ v bazi lastnih vektorjev, in dobimo
  \[
	\rho(x, A) = \frac{\sum_i \alpha_i^2 \lambda_i}{\sum_i \alpha_i^2}.
  \]
\end{proof}

\begin{izrek}[Courant-Fisher]
  Za simetrično matriko $A$ velja
  \[
	\lambda_i = \min_{\dim S = n-i+1} \max_{x \in S} \rho(x, A)
	= \max_{\dim R = i} \min_{x \in R} \rho(x, A).
  \]
\end{izrek}

\begin{proof}
  Označimo prvi izraz z $L$ in drugega z $D$.
  Za poljuben par podprostorov $S$ in $R$ je $\dim S + \dim R = n+1 > n$, torej
  obstaja neničelni vektor $x_{RS}$ v preseku.
  Naj bo na levi strani minimum dosežen pri $\hat{S}$, na desni pa pri
  $\hat{R}$.
  Potem je
  \[
	\max_{x \in \hat{S}} \rho(x, A) \ge \rho(x_{\hat{R} \hat{S}}, A) \ge \min_{x
	\in \hat{R}} \rho(x, A),
  \]
  iz česar sledi $L \ge D$.
  Če pa izberemo $\tilde{S} = \operatorname{Lin}(x_i, x_{i+1}, \ldots, x_n)$ in
  $\tilde{R} = \operatorname{Lin}(x_1, \ldots, x_i)$, pa je očitno
  \[
	L \le \max_{x \in \tilde{S}} \rho(x, A) = \lambda_i = \min_{x \in \tilde{R}}
	\rho(x, A) \le D,
  \]
  torej $L = D = \lambda_i$.
\end{proof}

\vprasanje{Povej in dokaži Courant-Fisherjev izrek.}

\begin{posledica}
  Če sta $A$ in $E$ simetrični matriki, potem za lastne vrednosti
  $\tilde{\lambda}_n \le \cdots \le \tilde{\lambda}_1$ matrike $A + E$ velja
  \[
	\lambda_i + \lambda_n(E) \le \tilde{\lambda}_i \le \lambda_i + \lambda_1(E).
  \]
\end{posledica}

\begin{proof}
  Uporabimo Courant-Fisherjev izrek za $\tilde{\lambda}_i$.
  Velja
  \[
	\tilde{\lambda}_i = \min_{\dim S = n-i+1} \max_{x \in S} \rho(x, A+E),
  \]
  upoštevamo $\rho(x, A+E) = \rho(x, A) + \rho(x, E)$ in ocenimo $\lambda_n(E)
  \le \rho(x, E) \le \lambda_1(E)$.
\end{proof}

\begin{posledica}[Weylov izrek]
  Če so $\lambda_n \le \cdots \le \lambda_1$ lastne vrednosti $A = A^T$ in
  $\tilde{\lambda}_n \le \cdots \le \tilde{\lambda}_1$ lastne vrednosti $A + E$
  za simetrično $E$, je $\abs{\tilde{\lambda}_i - \lambda_i} \le \norm{E}_2$.
\end{posledica}

\begin{proof}
  Velja $\norm{E}_2 = \max_i \abs{\lambda_i(E)}$.
\end{proof}

\vprasanje{Povej in dokaži Weylov izrek.}

\podnaslov{Rayleighova iteracija}

\begin{algorithm}
  \caption{Rayleighova iteracija}
  \label{alg:rayleighova-iteracija}
  \begin{algorithmic}
	\State Izberi $z_0 \ne 0$
	\For{$k=0, 1, 2, \ldots$}
	\State $\sigma_k = \rho(z_k, A)$
	\State Reši $(A - \sigma_k I) y_{k+1} = z_k$
	\State $z_{k+1} = y_{k+1} / \norm{y_{k+1}}_2$
	\EndFor
  \end{algorithmic}
\end{algorithm}

\begin{lema}
  Naj bo $\norm{z_0}_2 = 1$ približek za lastni vektor $x$ simetrične matrike
  $A$ z lastnimi vrednostmi $\abs{\lambda_n} \le \cdots \le \abs{\lambda_2} <
  \abs{\lambda_1}$.
  Če izvedemo en korak potenčne metode za $A$ z začetnim vektorjem $z_0$, potem
  za $z_1$ velja
  \[
	\norm{z_i \pm x_1} \le \frac{\abs{\lambda_2}}{\abs{\lambda_1}} \norm{z_0 -
	  x_1}_2 + O(\norm{z_0 - x_1}_2^2),
  \]
  kjer vzamemo tisto možnost v $z_1 \pm x_1$, ki da manjšo normo.
\end{lema}

\begin{proof}
  Predpostavimo, da je $z_0$ dober približek za $x_1$, torej
  \[
	z_0 = x_1 + \sum_{i=2}^n \alpha_i x_i
  \]
  za $\abs{\alpha_i} \ll 1$.
  Potem je
  \[
	z_1 \approx x_1 + \sum_{i=2}^n \alpha_i \frac{\lambda_i}{\lambda_1} x_i
  \]
  in
  \[
	\norm{z_1 - x_1}_2^2 = \sum_{i=2}^n \alpha_i^2
	\frac{\lambda_i^2}{\lambda_1^2}
	\le \frac{\lambda_2^2}{\lambda_1^2} \sum_{i=2}^n \alpha_i^2
	= \frac{\lambda_2^2}{\lambda_1^2} \norm{z_0 - x_1}^2.
  \]
\end{proof}

\vprasanje{Oceni napako po enem koraku potenčne metode.}

\begin{posledica}
  Naj bo $A$ simetrična in $\abs{\sigma - \lambda_i} \ll \abs{\sigma -
	\lambda_j}$ za $j \ne i$.
  Če je $z_0$ dober približek za lastni vektor $x_i$ in naredimo en korak
  inverzne iteracije, je
  \[
	\norm{z_1 \pm x_i} = O(\abs{\sigma - \lambda_i}) \norm{z_0 - x_i}_2
  \]
\end{posledica}

\begin{lema}
  Naj bo $A$ simetrična in $z$ približek za lastni vektor $x_i$ za $\lambda_i$.
  Potem je $\abs{\rho(z,A) - \lambda_i} \le 2 \norm{A}_2 \cdot \norm{z -
	x_i}_2^2$.
\end{lema}

\begin{proof}
  Brez škode za splošnost je $i = 1$.
  Naj bo $z = \sum_i \alpha_i x_i$ in $\sum_i \alpha_i^2 = 1$.
  Potem je
  \[
	\norm{z - x_1}_2^2 = (1 - \alpha_1)^2 + \alpha_2^2 + \cdots + \alpha_n^2
	= 2 - 2 \alpha_1.
  \]
  Velja $\rho(z, A) = \sum_i \lambda_i \alpha_i^2$, torej
  \[
	\rho(z, A) - \lambda_1
	= \sum_{i=1}^n \lambda_i \alpha_i^2 - \lambda_1 \sum_{i=1}^n \alpha_i^2
	= \sum_{i=2}^n (\lambda_i - \lambda_1) \alpha_i^2.
  \]
  Sledi
  \[
	\abs{\rho(z,A) - \lambda_1}
	\le \sum_{i=2}^n \abs{\lambda_i - \lambda_1} \alpha_i^2
	\le \sum_i (\abs{\lambda_i} + \abs{\lambda_1}) \alpha_1^2
	\le \sum_i 2 \norm{A}_2 \alpha_i^2
	\le 2 \norm{A}_2 (1 - \alpha_1^2)
  \]
  Predpostavimo, da je $\alpha_i \approx 1$, torej lahko to ocenimo še z
  \[
	2 \norm{A}_2 \cdot 2 (1 - \alpha_1)
	= 2 \norm{A}_2 \norm{z - x_1}_2^2.
  \]
\end{proof}

\vprasanje{Oceni napako, ki jo dobimo z Rayleighovim kvocientom.}

\begin{izrek}
  Rayleighova iteracija ima v bližini lastne vrednosti simetrične matrike
  kubično konvergenco.
\end{izrek}

\begin{proof}
  Naj bo $z_k$ blizu $x_i$.
  Vemo, da je
  \[
	\norm{z_{k+1} \pm x_i} = O(\abs{\sigma_k - \lambda_i} \norm{z_k - x_i}_2)
	= O(\norm{z_k - x_i}^2 \cdot \norm{z_k - x_i}_2).
  \]
\end{proof}

\vprasanje{Kakšen red konvergence ima Rayleighova iteracija za simetrične
  matrike? Dokaži.}

\podnaslov{QR iteracija}

Če je $A$ simetrična, začetna redukcija na zgornje Hessenbergovo obliko vrne
tridiagonalno simetrično matriko.
Tako $A$ reduciramo na
\[
  T =
  \begin{bmatrix}
	a_1 & b_1 & & & \\
	b_1 & a_2 & b_2 & & \\
		& b_2 & \ddots & \ddots & \\
		& & \ddots & a_{n-1} & b_{n-1} \\
	& & & b_{n-1} & a_n
  \end{bmatrix}.
\]
Algoritem bo še vedno porabil $O(n^3)$ operacij, ampak ga lahko zaradi simetrije
malo pohitrimo.
Predpostavimo lahko, da je $T$ nerazcepna.
En korak QR iteracije lahko izvedemo v $O(n)$ operacijah, če pa želimo
posodabljati še produkt matrik $Q_k$, je to dodatnih $O(n^2)$ operacij.

Imamo dva načina za izvajanje premika.
Prvi je Rayleighov premik $\sigma_k = \rho(e_n, T_k) = a_n^{(k)}$.

\begin{izrek}
  Če izvajamo QR iteracijo za tridiagonalno simetrično matriko $T$, potem se
  Rayleighovi premiki $\sigma_k$ ujemamo z Rayleighovimi kvocienti, ki bi jih
  dobili pri Rayleighovi iteraciji z $z_0 = e_n$.
\end{izrek}

\vprasanje{Kako deluje QR iteracija z Rayleighovimi premiki?}

Druga možnost so Wilkinsonovi premiki, kjer pogledamo
\[
  W_k =
  \begin{bmatrix}
	a_{n-1}^{(k)} & b_{n-1}^{(k)} \\
	b_{n-1}^{(k)} & a_n^{(k)}
  \end{bmatrix}
\]
in izberemo tisto lastno vrednost, ki je bližja $a_n^{(k)}$.
Tukaj imamo celo globalno konvergenco, ki je v bližini lastne vrednosti kubična.

\vprasanje{Kaj so Wilkinsonovi premiki?}

\podnaslov{Bisekcija}

\begin{izrek}
  Če je $T$ nerazcepna simetrična tridiagonalna matrika, ima vse lastne
  vrednosti enostavne.
\end{izrek}

\begin{proof}
  Naj bo $\lambda$ lastna vrednost $T$.
  Ker je matrika simetrična, sta geometrijska in algebraična večkratnost enaki.
  Ker so $b_i \ne 0$, je prvih $n-1$ stolpcev $T - \lambda I$ neodvisnih, torej
  je rang te matrike enak $n-1$ in je večkratnost $\lambda$ enaka $1$.
\end{proof}

\vprasanje{Dokaži: nerazcepna simetrična tridiagonalna matrika ima vse lastne
  vrednosti enostavne.}

\begin{definicija}
  Matriki $A$ in $B$ sta \pojem{kongruentni}, če obstaja nesingularna $X$, da je
  $B = X^T A X$.
\end{definicija}

Če je $A$ simetrična, je $\varphi(x) = x^T A x$ kvadratna forma.
Potem za $x = Zy$ dobimo $x^T A x = y^T Z^T A Z y$ v drugi bazi.

\begin{definicija}
  Za vsako $n \times n$ simetrično matriko $A$ definiramo \pojem{inercijo} $A$
  kot trojico $(\nu, z, p)$, kjer je $\nu$ število negativnih, $z$ število
  ničelnih in $p$ število pozitivnih lastnih vrednosti.
\end{definicija}

\begin{izrek}[Sylvester]
  Kongruentni simetrični matriki imata enako inercijo.
\end{izrek}

\begin{proof}
  Naj bo $(\nu, z, p)$ inercija $A$ in $(\nu', z', p')$ inercija $X^T A X$.
  Recimo $\nu' < \nu$.
  Vsi lastni vektorji matrike $A$ za negativne lastne vrednosti razpenjajo nek
  invarianten podprostor $\mathcal{N}$ dimenzije $\nu$.
  Podobno vsi lastni vektorji nenegativnih lastnih vrednosti $X^T A X$
  razpenjajo invarianten podprostor $\mathcal{P}$ dimenzije $n - \nu'$ za $X^T A
  X$.
  Prostor $X \mathcal{P}$ je tudi podprostor dimenzije $n - \nu'$, ki se mora
  sekati z $\mathcal{N}$ po predpostavki $\nu' < \nu$.
  Naj bo $v \in \mathcal{N} \presek X \mathcal{P}$ neničeln.
  Velja
  \[
	\rho(v, A) = \frac{v^T A v}{v^T v} < 0,
  \]
  po drugi strani pa obstaja $w \in \mathcal{P}$, da je $v = X w$.
  Velja
  \[
	\rho(v, A) = \frac{w^T X^T A X w}{w^T X^T X w} \ge 0,
  \]
  kar je protislovje.
  Podobno v primeru $p' < p$.
\end{proof}

\vprasanje{Povej in dokaži Sylvestrov izrek.}

Kako izračunamo inercijo za $T - \lambda I$?
Uporabimo $L D L^T$ razcep, z
\[
  L =
  \begin{bmatrix}
	1 & & & & \\
	l_1 & 1 & & & \\
	  & l_2 & 1 & & \\
	  & & \ddots & \ddots & \\
	& & & l_{n-1} & 1
  \end{bmatrix}
\]
in $D = \operatorname{diag}(d_1, \ldots, d_n)$.
Razcep poiščemo podobno kot LU razcep.
Vemo, da je inercija $T - \lambda I$ enaka inerciji $D$.
Lastne vrednosti poiščemo z bisekcijo, kjer upoštevamo informacije iz inercije.

V postopku za račun $LDL^T$ razcepa smo nekje delili z diagonalnim elementom,
kar lahko načeloma da vmesni rezultat $d_i = -\infty$.
To nas ne moti, tako vrednost štejemo kot negativno, v naslednjem koraku pa bomo
tako in tako dobili $0$.

\vprasanje{Opiši metodo bisekcije za iskanje lastnih vrednosti.}

\podnaslov{Jacobijeva metoda}

V Jacobijevi metodi matrike ne reduciramo na tridiagonalno obliko.
Vzamemo Jacobijevo (Givensovo) rotacijo
\[
  R_{pq}^T =
  \begin{bmatrix}
	1 &        &   &    &        &   &        & \\
	  & \ddots &   &    &        &   &        & \\
	  &        & 1 &    &        &   &        & \\
	  &        &   & c  &        & s &        & \\
	  &        &   &    & \ddots &   &        & \\
	  &        &   & -s &        & c &        & \\
	  &        &   &    &        &   & \ddots & \\
	  &        &   &    &        &   &        & 1
  \end{bmatrix},
\]
kjer izberemo $c$ in $s$ tako, da bo $\tilde{a}_{pq} = \tilde{a}_{qp} = 0$.
S tem si lahko pokvarimo prejšnja postavljanja.

\begin{lema}
  Če $\tilde{A}$ dobimo iz $A$ kot $\tilde{A} = R_{pq}^T A R_{pq}$, kjer
  $R_{pq}$ določimo tako, da je $\tilde{a}_{pq} = \tilde{a}_{qp} = 0$, potem
  $\operatorname{off}(\tilde{A})^2 = \operatorname{off}(A)^2 - 2 a_{pq}^2$,
  kjer je $\operatorname{off}$ norma izvendiagonalnega dela.
\end{lema}

\begin{proof}
  Ker je $R_{pq}$ ortogonalna, velja
  \[
	\operatorname{off}(\tilde{A})^2 + \sum_i \tilde{a}_{ii}^2
	= \norm{\tilde{A}}_F^2
	= \norm{A}_F^2
	= \operatorname{off}(A)^2 + \sum_i a_{ii}^2.
  \]
  Vsi diagonalni elementi razen $pp$ in $qq$ se ujemajo, torej se to poenostavi
  v
  \[
	\operatorname{off}(\tilde{A})^2 = \operatorname{off}(A)^2 + a_{pp}^2 +
	a_{qq}^2 - \tilde{a}_{pp}^2 - \tilde{a}_{qq}^2.
  \]
  Tudi $2\times 2$ podmatriki mest $(p, q)$ imata enako Frobeniusovo normo,
  torej je ostanek enak $- 2 a_{pq}^2$.
\end{proof}

\vprasanje{Dokaži, da se pri Jacobijevi iteraciji norma izvendiagonalnih
  elementov zmanjšuje z iteracijami.}

S preprostim izračunom lahko pridemo do naslednjega predpisa za $s$ in $c$:
\begin{gather*}
  \tau = \frac{a_{pp} - a_{qq}}{2 a_{pq}}, \\
  t = \frac{\operatorname{sgn} \tau}{\abs{\tau} + \sqrt{1 + \tau^2}}, \\
  c = \inv{\sqrt{1 + t^2}}, \\
  s = ct.
\end{gather*}
V enem koraku metode potem posodobimo $A = R_{pq}^T A R_{pq}$ in $Q = Q R_{pq}$.

\vprasanje{Opiši en korak Jacobijeve iteracije.}

Jacobijeva metoda je običajno počasnejša od ostalih, lahko pa z njo natančneje
izračunamo majhne lastne vrednosti spd matrik.
Metoda ima več variant:
\begin{itemize}
\item \pojem{Klasična varianta}: v vsakem koraku uničimo po absolutni vrednosti
  največji izvendiagonalni element.
  Preverjanje, kateri element je največji, lahko optimiziramo iz $O(n^2)$ v
  $O(n)$, če si shranjujemo največji element po vrsticah in stolpcih.
\item \pojem{Ciklična varianta}: elemente uničujemo po vrsti, na koncu se vrnemo
  na začetek.
\item \pojem{Pragovna varianta}: podobno kot ciklična metoda, a uničimo le
  dovolj velike elemente.
\end{itemize}

\vprasanje{Opiši tri variante Jacobijeve metode.}

\podnaslov{Deli in vladaj}

Naj bo $T$ nerazcepna simetrična tridiagonalna $n \times n$ matrika in $m
\approx \frac{n}{2}$.
Razdelimo $T$ na dve matriki;
\[
  T =
  \begin{bmatrix}
	T_1 & \\
	& T_2
  \end{bmatrix}
  + b_m v v^T,
\]
kjer je $v = e_m + e_{m+1}$, matriki $T_1$ in $T_2$ pa sta nerazcepni
tridiagonalni.
Rekurzivno lahko poiščemo $T_1 = Q_1 D_1 Q_1^T$ in $T_2 = Q_2 D_2 Q_2^T$.
Potem velja
\[
  T =
  \begin{bmatrix}
	Q_1 & \\
	& Q_2
  \end{bmatrix}
  \left(
	\begin{bmatrix}
	  D_1 & \\
		  & D_2
	\end{bmatrix}
	+ b_m u u^T
  \right)
  \begin{bmatrix}
	Q_1^T & \\
	& Q_2^T
  \end{bmatrix}
\]
za
\[
  u =
  \begin{bmatrix}
	Q_1^T & \\
		  & Q_2^T
  \end{bmatrix}
  v.
\]
Recimo, da znamo učinkovito rešiti problem lastnih vrednosti za matriko $D +
\rho u u^T = \tilde{Q} \Lambda \tilde{Q}^T$.
Potem je
\[
  T =
  \begin{bmatrix}
	Q_1 & \\
	& Q_2
  \end{bmatrix}
  \tilde{Q}
  \Lambda
  \tilde{Q}^T
  \begin{bmatrix}
	Q_1^T & \\
	& Q_2^T
  \end{bmatrix}.
\]

\begin{izrek}
  Naj bo $A = D + \rho u u^T$ za $D = \operatorname{diag}(d_1, \ldots, d_n)$ in
  $d_1 \ge d_2 \ge \cdots \ge d_n$.
  \begin{itemize}
  \item Če je $u_i = 0$, je $d_i$ lastna vrednost $A$ z lastnim vektorjem $e_i$.
  \item Če je $d_i = d_{i+1}$, potem je $d_i$ lastna vrednost $A$ za lastni
	vektor $- u_{i+1} e_i + u_i e_{i+1}$.
  \end{itemize}
\end{izrek}

\vprasanje{Opiši osnovno delovanje metode deli in vladaj. Kakšni so posebni
  primeri za lastne vrednosti?}

Preostane problem lastnih vrednosti za $A = D + \rho u u^T$, kjer je $d_1 >
\cdots > d_n$ in $u_i \ne 0$ za vsak $i$.
Naj bo $\lambda$ lastna vrednost $A$.
Predpostavimo, da je $\lambda \ne d_i$.
Vemo, da je $A - \lambda I$ singularna.
Ker je
\[
  A - \lambda I = D - \lambda I + \rho u u^T
  = (D - \lambda I)(I + \rho (D - \lambda I)^{-1} u u^T)
\]
in ker je prva matrika v produktu nesingularna, je druga singularna.
Za $x = (D - \lambda I)^{-1} u$ in $y^T = u^T$ velja naslednja lema:

\begin{lema}
  $\det (I + x y^T) = 1 + y^T x$.
\end{lema}

\begin{proof}
  Če je $z \bot y$, je $(I + xy^T)z = z$, takih je $n-1$ linearno neodvisnih
  $z$, torej je $1$ $(n-1)$-kratna lastna vrednost.
  Velja $(I + xy^T)x = (1 + y^T x)x$, torej je $1 + y^T x$ zadnja lastna
  vrednost.
\end{proof}

Računamo
\[
  \det (I+ \rho (D - \lambda I)^{-1} u u^T) = 1 + \rho u^T (D - \lambda I)^{-1}
  u
  = 1 + \rho \sum_{i=1}^n \frac{u_i^2}{d_i - \lambda} = f(\lambda),
\]
čemur pravimo \pojem{sekularna funkcija}.

\vprasanje{Izpelji sekularno funkcijo.}

Lastne vrednosti se rešitve sekularne enačbe $f(\lambda) = 0$.
Funkcija ima $n$ enostavnih polov, za $\rho > 0$ je naraščajoča in za $\rho < 0$
padajoča z asimptoto $y = 1$.
Torej ima natanko $n$ ničel na znanih intervalih $(d_i, d_{i-1})$ oziroma na
enem od neskončnih intervalov.
Če je $\alpha$ lastna vrednost, lahko izračunamo lastni vektor $(D - \alpha
I)^{-1} u$.

Krivulja je strma okoli ničle in položna na velikem delu intervala, kar je
težava, ker tako tangentna metoda kot bisekcija ne delujeta dobro.
Namesto tega $f$ aproksimiramo z
\[
  h(x) = \frac{c_1}{d_{i+1} - \lambda} + \frac{c_2}{d_i - \lambda} + c_3,
\]
kjer konstante izračunamo tako, da se vrednosti $h$ in $f$ ter njunih odvodov
ujemajo v približku $x_0$.

\vprasanje{Kako poiščeš lastne vrednosti in lastne vektorje matrike $D + \rho u
  u^T$?}

Metoda je hitrejša od QR iteracije za $n > 30$, če potrebujemo tudi lastne
vektorje.
Ne deluje pa brez računanja $Q$.

% LocalWords:  Courant-Fisherjev Courant-Fisher Weylov Rayleighova Rayleighovi
% LocalWords:  Rayleighovimi Wilkinsonovi Sylvester invarianten LU Jacobijevi
% LocalWords:  Givensovo izvendiagonalnega izvendiagonalnih izvendiagonalni spd
% LocalWords:  Pragovna tridiagonalni
